\label{sec:methodology}

\subsection{Algorithmic Map Generation}

We generate \textsc{bisect} congressional maps for NC ($k=14$), WI ($k=8$), and TX
($k=38$) using 2020 Census PL 94-171 population data and TIGER/Line tract geometries.
The algorithm is invoked as:
\begin{verbatim}
bisect build <state>_incumb_test --year 2020
    --structure prime-factor
    --weights-override geographic
    --workers 8
\end{verbatim}
The \texttt{--weights-override geographic} flag instructs the algorithm to use only
geographic edge weights (shared boundary length between adjacent tracts), without
incorporating any precinct-level partisan data. This ensures that the algorithm has
no information about where Democratic or Republican voters are concentrated.

Single-seed results use seed\,$=42$ and are marked with $^\dagger$.

\subsection{Partisan Data Assignment}

After generating the \textsc{bisect} district assignments, we assign 2020 VEST
presidential precinct returns to each district by:
\begin{enumerate}
  \item Performing an area-weighted spatial join of VEST precincts to 2020 census
    tract polygons, producing tract-level estimates of $V_t^D$ and $V_t^R$.
  \item Summing $V_t^D$ and $V_t^R$ across all tracts $t$ assigned to each district $d$
    to produce district-level vote totals.
  \item Computing $v_d^D = V_d^D / (V_d^D + V_d^R)$.
\end{enumerate}

The same procedure is applied to the enacted 2022 congressional district boundaries.
This ensures that any differences in safe-seat counts between algorithmic and enacted
maps are attributable to differences in district boundary placement, not to differences
in how partisan data is assigned.

\subsection{Competitive District Count}

A district is classified as competitive if $|v_d^D - 0.50| \leq 0.05$ (within 5
percentage points of parity), lean-D if $v_d^D \in (0.50, 0.65]$, lean-R if
$v_d^D \in [0.35, 0.50)$, safe-D if $v_d^D > 0.65$, and safe-R if $v_d^D < 0.35$.
Each district falls into exactly one category.

\subsection{Statistical Testing}

To test whether the difference in safe-seat fractions between \textsc{bisect} and
enacted maps is statistically significant, we use a two-sample test of proportions.
Let $f_\text{alg}$ be the algorithmic safe-seat fraction and $f_\text{enact}$ be
the enacted safe-seat fraction, with sample sizes $k_\text{alg} = k_\text{enact} = k$.
The test statistic is:
\[
  z = \frac{f_\text{enact} - f_\text{alg}}{\sqrt{\hat{p}(1-\hat{p})(1/k + 1/k)}},
\]
where $\hat{p} = (f_\text{enact} + f_\text{alg})/2$ is the pooled proportion.
We report the one-sided $p$-value for the hypothesis that $f_\text{enact} > f_\text{alg}$.

Note that the single-seed nature of the \textsc{bisect} results introduces additional
uncertainty not captured by the binomial test; the reported $p$-values should be
interpreted as lower bounds on the evidence against the null.
